\chapter{Data-Structure Attribution: Naming Every Byte}
\label{ch:attribution}

A kernel-level claim (“this kernel is memory-bound”) is not yet an architect’s claim (“the keyframe database belongs near storage”). Bridging that gap requires knowing, for each byte of DRAM traffic, \emph{which named data structure it belongs to}. This chapter presents the TaggedAllocator attribution: the instrument (\S\ref{sec:attr-design}), the coverage discipline (\S\ref{sec:attr-coverage}), the static-budget finding (\S\ref{sec:attr-static}), the per-kernel attribution results (\S\ref{sec:attr-results}), the register-spill finding (\S\ref{sec:attr-spill}), and the bounded residuals (\S\ref{sec:attr-residuals}).

\section{The Three-Layer Instrument}
\label{sec:attr-design}

Attribution is a chain of fragile assumptions (address $\rightarrow$ allocation $\rightarrow$ owning code $\rightarrow$ data-structure name), so the design uses three independent layers for cross-checking:

 \textbf{Layer 1: source-side journal (semantic layer).} Every GPU allocation in cuVSLAM flows through allocation wrappers; an injected journal records, for each allocation and free, the pointer, size, wrapper kind, and a host \emph{call stack}. Resolved offline against debug symbols, the innermost non-plumbing frame names the owning data structure (e.g., the Schur-complement solver’s constructor $\Rightarrow$ \code{ba\_linear\_system}). A fixed tag vocabulary (\code{images\_raw}, \code{pyramid\_levels}, etc.) keeps names stable across runs and sequences. Pinned host buffers (device-visible under unified addressing) are journaled with a distinguishing suffix, as their traffic crosses PCIe rather than DRAM.

\textbf{Layer 2: driver-side sidecar (temporal layer).} The NVBit tool independently logs every driver-level allocation and free, stamped with the current kernel-launch counter, the same clock the address trace uses. This provides allocation lifetimes and catches allocations that bypass wrappers (e.g., math libraries’ scratch), which are tagged as driver-internal rather than silently mislabeled.

\textbf{Layer 3: the join.} A streaming pass over the address trace maintains the live-allocation set in launch order and resolves each global-space access to its containing allocation’s tag; shared-space and local-space records are bucketed as on-chip and spill, respectively (\S\ref{sec:loc-correction}). Anything unresolved lands in an explicit \code{unmapped} bucket, the join’s own error bar, reported and bounded.

\section{Coverage: Auditing the Windows}
\label{sec:attr-coverage}

Traces are captured in windows, and sparse kernels can drift out of a window between runs (launch indices shift with keyframe timing). The pipeline treats coverage as a proof obligation: after each campaign, an audit diffs every sequence’s \emph{full} launch map (from a cheap uninstrumented pass) against the union of attributed kernels, and a planner captures gap-fill windows anchored on dense clusters of missing kernels’ launches (clusters are robust; isolated occurrences drift). Across the 27-sequence campaign this loop converged to \textbf{zero missing kernels}, with all audit artifacts committed. Without this discipline, three sequences would have silently lacked the loop-closure scan and one the keyframe-refresh cluster, precisely the architecturally interesting kernels.

\section{The GPU Memory Budget Is Static}
\label{sec:attr-static}

The journal’s first result is structural. Over a full 2{,}500-frame loop-closing session, cuVSLAM performs \textbf{240 device allocations totalling 108.65\,MB all at initialization, none freed mid-run and the identical 240 in a 30-frame run}. The budget by tag (Table~\ref{tab:budget}): images and pyramids $\sim$71\,MB; the BA linear system 24.3\,MB (plus a 15.4\,MB pinned-host mirror); and most decisively a \textbf{fixed 6.7\,MB keyframe-descriptor buffer}. The keyframe \emph{database} (the only data structure that grows with session length) does not grow on the GPU at all: its growing store lives host-side in the memory-mapped database, confirmed by direct host measurement (708\,MB peak host memory against the static GPU budget; \S\ref{sec:syn-host}). The in-storage verdict of Chapter~\ref{ch:synthesis} rests on this measured division of labour, not on inference from kernel behaviour.

\begin{table}[t]
\centering
\caption{The static GPU allocation budget (full session; peak = total, nothing freed mid-run).}
\label{tab:budget}
\begin{tabular}{@{}lrr@{}}
\toprule
tag & allocations & peak live (MB) \\
\midrule
\code{images\_raw} & 92 & 39.97 \\
\code{pyramid\_levels} & 99 & 31.15 \\
\code{ba\_linear\_system} & 31 & 24.26 \\
\code{keyframe\_descriptors} & 2 & \textbf{6.73 (fixed)} \\
\code{feature\_tracks} & 8 & 4.27 \\
\code{depth\_pyramid} & 3 & 2.21 \\
\code{icp\_state} & 4 & 0.07 \\
\midrule
device total & 240 & 108.65 \\
pinned host (device-visible) & 34 & 17.14 \\
\bottomrule
\end{tabular}
\end{table}

\section{Attribution Results}
\label{sec:attr-results}

The join resolves \textbf{274 of 274} journaled allocations (240 device + 34 pinned-host) with zero unknowns. Across the full 27-sequence campaign, \textbf{48 of 49 kernels receive a unanimous top data-structure tag in every sequence they appear in} the attribution analogue of the 91\% class consistency of \S\ref{sec:char-stability}, and the license to make per-kernel data-structure claims input-independently. Selected rows (steady-state window; shares are sector-weighted):

\begin{itemize}
  \item \stTrack{} (loop-closure scan): \textbf{92.6\% register spill}, with the global remainder on \code{keyframe\_descriptors}; unmapped $<$1\%.
  \item \stBuild{} (keyframe ingest): 93\% \code{keyframe\_descriptors}.
  \item \code{sba::reduced\_system\_stage\_2} (solver core): \textbf{96.9\% \code{ba\_linear\_system}} one named, pre-sized, contiguous structure carrying essentially all of the solver’s DRAM traffic.
  \item Front-end compute kernels (convolutions, feature selection, sorting): 83--98\% of accesses are shared-memory tiles (on-chip; no DRAM claim), with the global residue falling on the pyramid, image, and track structures the pipeline stage implies.
\end{itemize}

Mode differences surface exactly where expected: stereo configurations journal no depth or ICP structures; loop-free KITTI sequences show \stTrack{} launches only when their trajectories genuinely close loops the attribution reproduces workload semantics, itself a correctness check.

\section{The Register-Spill Finding}
\label{sec:attr-spill}

The single most consequential attribution row is \stTrack{}'s (Figure~\ref{fig:attrib}): the kernel motivating the near-data discussion spends \textbf{92.6\% of its DRAM volume on register spill} compiler-generated scratch traffic, invisible to any counter-only methodology.

\begin{figure}[t]
\centering
\begin{tikzpicture}
\begin{axis}[
  width=0.82\textwidth, height=4.4cm,
  xbar stacked, xmin=0, xmax=100,
  xlabel={share of kernel DRAM traffic (\%)}, xlabel style={font=\small},
  ytick=data, symbolic y coords={sba::reduced\_2,st\_build\_cache,st\_track},
  yticklabel style={font=\footnotesize\ttfamily}, tick label style={font=\scriptsize},
  legend style={at={(0.5,-0.42)},anchor=north,legend columns=3,font=\scriptsize,draw=black!30},
  bar width=13pt, enlarge y limits=0.35,
]
\addplot[fill=cPIM!55,draw=black!40] coordinates {(92.6,st\_track)(4,st\_build\_cache)(2,sba::reduced\_2)};
\addplot[fill=black!12,draw=black!40] coordinates {(0.6,st\_track)(3,st\_build\_cache)(1.1,sba::reduced\_2)};
\addplot[fill=cISP!45,draw=black!40] coordinates {(6.8,st\_track)(93,st\_build\_cache)(96.9,sba::reduced\_2)};
\legend{register spill (scratch),unmapped/other,named global data structure}
\end{axis}
\end{tikzpicture}
\caption[Per-kernel memory-space attribution.]{Memory-space attribution for three
representative kernels. The loop-closure scan \stTrack{} spends 92.6\% of its DRAM
volume on register spill (a compiler artifact, not data) while its true data
gather is the small remainder on \code{keyframe\_descriptors}. By contrast the BA
solver core and the keyframe-ingest kernel spend almost all of their traffic on a
single named structure.}
\label{fig:attrib}
\end{figure} The finding cuts both ways. For \emph{this} workload: the loop-closure scan’s bandwidth bill is mostly a register-pressure artifact (per-thread descriptor state exceeds register budget), so a register-file/occupancy trade or spill-aware compilation attacks the volume, while the residual the true database gather (\S\ref{sec:loc-correction}) is the scatter-PiM candidate. For the \emph{methodology}: spill is a first-class DRAM citizen that only per-space attribution can expose; a caveat applies (register allocation is a codegen decision for this compiler/target; taxonomy transfers, spill \emph{quantity} must be re-derived per target, and cross-target validation is scheduled work).

\section{Bounded Residuals}
\label{sec:attr-residuals}

The join’s honesty budget: for most kernels, the unmapped share is below 1--7\%. Two systematic residuals remain, both bounded and explained. First, one solver-assembly kernel carries $\sim$40\% unmapped traffic \emph{identically across all 27 sequences} the signature of static module memory (device-side globals allocated by the module loader, invisible to both journal layers), not a lurking data structure; analysis names the two kernels (\code{sba::build\_full\_system}, \code{lk\_track}) that carry 83\% of all unmapped traffic, making the residual a prioritized target for a planned kernel-argument correlation layer. Second, texture-path reads (the LK tracker samples pyramids through texture units) never appear in the address trace, so image-structure traffic is a stated lower bound, with counter-side texture totals corroborating the affected kernels. Neither residual moves any verdict: affected kernels' classes rest on counters and stage evidence, independent of the unmapped share.